% =====================================================================
% Springer Nature jurnal maqolasi shabloni (sn-jnl sinfi) - Ilmiy Yo'l
% Izohlar o'zbek tilida, maqola matni inglizcha yoziladi.
%
% DIQQAT: sn-jnl.cls va sn-*.bst fayllari Springer Nature'ning rasmiy
% fayllari. Ularni shu papkaga qo'ying (Overleaf'dagi "Springer Nature
% LaTeX Template" ichida tayyor bor) yoki Springer'ning mualliflar
% sahifasidan oling.
%
% Kompilyatsiya: pdflatex -> bibtex -> pdflatex -> pdflatex
%
% MUHIM: jurnalning "Instructions for Authors" sahifasini o'qing.
% Havolalar uslubini sinf parametri belgilaydi (pastda).
% =====================================================================

% Havolalar uslubi: sn-mathphys-num (raqamli [1]), sn-basic, sn-vancouver,
% sn-apacite (muallif-yil) va boshqalar. Jurnal talabiga qarab tanlang.
\documentclass[sn-mathphys-num]{sn-jnl}

\usepackage{graphicx}
\usepackage{multirow}
\usepackage{amsmath,amssymb,amsfonts}
\usepackage{xcolor}
\usepackage{textcomp}
\usepackage{booktabs}

\raggedbottom

\begin{document}

% --- Sarlavha: aniq, mazmunli, qisqa. Qavs ichida ixcham sarlavha. ---
\title[Short Title]{Title of Your Article: Specific, Informative and Concise}

% --- Mualliflar: yulduzcha (*) muxbir muallifni bildiradi ---
\author*[1]{\fnm{First} \sur{Author}}\email{first.author@example.com}

\author[1]{\fnm{Second} \sur{Author}}\email{second.author@example.com}

\author[2]{\fnm{Third} \sur{Author}}\email{third.author@example.com}

\affil*[1]{\orgdiv{Department of Information Technologies}, \orgname{Example University}, \orgaddress{\city{City}, \postcode{100000}, \country{Country}}}

\affil[2]{\orgdiv{Department}, \orgname{Example Institute}, \orgaddress{\city{City}, \postcode{100000}, \country{Country}}}

% --- Abstract: muammo, usul, asosiy natija, xulosa. Bitta xatboshi. ---
\abstract{State the problem in one sentence. Describe the method you propose or the data you analyse. Report the main quantitative result with a number. End with the practical or scientific significance. Keep the abstract self-contained, without citations. Check the journal's word limit.}

\keywords{keyword one, keyword two, keyword three, keyword four}

\maketitle

% =====================================================================
% Introduction: muammo, dolzarbligi, oldingi ishlar, research gap, hissa.
% =====================================================================
\section{Introduction}\label{sec1}
State the research problem and why it matters to the field \cite{he2016resnet,redmon2016yolo}. Summarise what earlier studies have done and where they fall short. Name the gap clearly: new data, a new method, or a new context.

State the contribution of this article in one sentence. Finish with a short outline of the remaining sections.

% =====================================================================
% Methods: takrorlash mumkin bo'ladigan darajada aniq.
% =====================================================================
\section{Methods}\label{sec2}
Describe the data set, its source and size. Describe the proposed method step by step. Give the evaluation metric as a formula, for example the $F_1$ score:
\begin{equation}
  F_1 = \frac{2PR}{P+R},
  \label{eq:f1}
\end{equation}
where $P$ is precision and $R$ is recall. Refer to equations as Eq.~(\ref{eq:f1}).

% =====================================================================
% Results: faqat natijalar, izohsiz.
% =====================================================================
\section{Results}\label{sec3}
Present the results without interpretation. Table~\ref{tab:results} compares the proposed method with the baseline. Replace the placeholders with your own measured values.

\begin{table}[h]
  \caption{Comparison with the baseline (replace with your results)}
  \label{tab:results}
  \begin{tabular}{@{}lcc@{}}
    \toprule
    Method & $F_1$ score (\%) & Time (ms) \\
    \midrule
    Baseline & XX.X & XX \\
    Proposed & XX.X & XX \\
    \botrule
  \end{tabular}
\end{table}

% Rasm qo'shish uchun (fayl papkada bo'lishi kerak, masalan fig1.pdf):
% \begin{figure}[h]
%   \centering
%   \includegraphics[width=0.9\textwidth]{fig1.pdf}
%   \caption{Short, self-explanatory caption.}
%   \label{fig:overview}
% \end{figure}

% =====================================================================
% Discussion: natijalar nimani anglatadi, cheklovlar, boshqa ishlar bilan solishtirish.
% =====================================================================
\section{Discussion}\label{sec4}
Interpret the results. Explain why the proposed method behaves as observed, compare with related work \cite{vaswani2017attention} and state the limitations honestly.

\section{Conclusion}\label{sec5}
Summarise the findings in two or three sentences and state the future work. Do not repeat the abstract and do not introduce new results.

\bmhead{Acknowledgements}
Thank the people who supported the work. Write ``None'' if there are none.

% --- Springer ko'pincha quyidagi bandlarni talab qiladi. Jurnal talabiga qarang. ---
\section*{Declarations}
\begin{itemize}
\item Funding: state the funding source or ``No funding''.
\item Conflict of interest: state or ``The authors declare no competing interests''.
\item Data availability: where the data can be obtained.
\item Code availability: where the code can be obtained, if applicable.
\item Author contributions: who did what.
\end{itemize}

% --- Adabiyotlar: references.bib faylidan ---
\bibliography{references}

\end{document}
